\scanpage{8}
paraît (\unsure{Laborde|Lazard} dixit) que (sans condition que
$2$ soit inversible) les modules $(M,q)\oplus(M,-q)$
sont de la forme \inserted{$\alpha(E)=$} $E\oplus\check E$. (Mais ce n'est pas
\deleted{Dans}
vrai dans \deleted{\unsure{tout}} \inserted{le} cas d'un topos loc. annelé —
p.ex. cela n'est pas vrai pour le topos classifiant
$B_{O(n)}$ si $n\ge1$)

\margin{à vérifier}
Dans l'un et l'autre cas, $W(S)$ peut
dès lors se construire directement comme
quotient de l'ens des classes d'isom.
de modules quadratiques, avec
$(M,q)\sim(M',q')$ si \deleted{$(M,q)\oplus(M',-q')$}
\deleted{est isom. à un $E\oplus\check E$} \inserted{$=\alpha(E)$}
$\exists$ Mod loc. lib. de r.f. $E,F$ avec
\[
(M,q)\oplus(M',-q')+\alpha(F)\simeq\alpha(E)
\]

[\uline{NB} Si $F$ est \unsure{isom.} à un fact direct de $E$,
\inserted{et si on dispose d'un ``th de Witt''}
— ce qui est tjrs le cas si $S$ est le spectre d'
un corps — \deleted{et si} alors cette condition signifie
aussi que $\exists G$ loc. lib de r.f. tel que
\[
(M,q)\oplus(M',-q')\simeq\alpha(G)
\]
]

On voit que $\alpha(K(S))$ est un \uline{idéal} dans
$WG(S)$ : donc $W(S)$ est encore un anneau.
Mais il n'est plus augmenté vers $H^0(S,\mathbb Z)$,
seulement vers $H^0(S,\underbrace{\mathbb Z/2\mathbb Z}_{\mathbb F_2})$ (donc vers
$\mathbb F_2$ si $S$ connexe $\ne\varnothing$). Mais que peut-on
dire de la pré-$\lambda$-structure de $WG(S)$
— passerait-elle au quotient ?

\mspage{7}
